% BEGIN VOL08-EXPANSION VIII05
\section{Supplementary worked examples}
\label{sec:viii05-pedagogy}

\begin{example}[Attaching a 1-cell]\label{ex:viii05-ped-01}
Attaching a 1-cell to two points by its boundary $S^0$ creates an interval joining them; attaching both endpoints to one point creates a circle.  The attaching map records the topology created at the boundary.
\end{example}

\begin{example}[Attaching a 2-cell to a circle]\label{ex:viii05-ped-02}
Attach $D^2$ to $S^1$ by the identity map on $S^1$.  The result is homeomorphic to $D^2\cup_{S^1}D^2\cong S^2$, exhibiting the sphere as a $0$-cell plus a $2$-cell.
\end{example}

\begin{example}[Degree-$m$ 2-cell attachment]\label{ex:viii05-ped-03}
Attaching a 2-cell to $S^1$ by a map of degree $m$ produces a standard two-dimensional CW complex whose first homology becomes $\mathbb Z/m$ for $m\ne0$.  The attaching degree becomes a cellular boundary coefficient.
\end{example}

\section{Graded supplementary exercises}
The following exercises progress from direct computation and construction to proof, hypothesis testing, application, and synthesis.

\subsection*{Standard computations and constructions}

\begin{exercise}[Boundary of a cell]\label{exr:viii05-ped-01}
What is the boundary sphere used to attach an $n$-cell?
\end{exercise}
\begin{hint}
An $n$-cell is a copy of $D^n$.
\end{hint}
\begin{solution}
Its boundary is $S^{n-1}=\partial D^n$.
\end{solution}

\begin{exercise}[Attach a 0-cell]\label{exr:viii05-ped-02}
What data are needed to attach a 0-cell?
\end{exercise}
\begin{hint}
The boundary of $D^0$ is empty.
\end{hint}
\begin{solution}
No attaching map data are needed; a 0-cell simply adds a point.
\end{solution}

\begin{exercise}[Circle from one cell]\label{exr:viii05-ped-03}
Describe $S^1$ as a CW complex with one 0-cell and one 1-cell.
\end{exercise}
\begin{hint}
Identify both endpoints of $D^1$ with the 0-cell.
\end{hint}
\begin{solution}
Attach the boundary $S^0=\{-1,1\}$ of the 1-cell to the same 0-cell; the quotient is a circle.
\end{solution}

\begin{exercise}[Sphere cells]\label{exr:viii05-ped-04}
Give the two-cell CW structure on $S^n$.
\end{exercise}
\begin{hint}
Use one 0-cell and one $n$-cell.
\end{hint}
\begin{solution}
Attach $D^n$ to a point by the unique map $S^{n-1}\to\{*\}$; the quotient $D^n/S^{n-1}$ is $S^n$.
\end{solution}

\begin{exercise}[Degree attachment]\label{exr:viii05-ped-05}
If a 2-cell is attached to $S^1$ by degree $m$, what integer will appear in the cellular boundary?
\end{exercise}
\begin{hint}
The boundary coefficient is the attaching degree.
\end{hint}
\begin{solution}
The differential $C_2\to C_1$ is multiplication by $m$ in the one-cell model.
\end{solution}

\subsection*{Proofs}

\begin{exercise}[Quotient description]\label{exr:viii05-ped-06}
Prove cell attachment $X\cup_f D^n$ is a pushout/quotient construction.
\end{exercise}
\begin{hint}
Identify each $u\in S^{n-1}$ with $f(u)\in X$.
\end{hint}
\begin{solution}
Take the disjoint union $X\sqcup D^n$ and impose precisely the relations $u\sim f(u)$ for $u\in\partial D^n$; this realizes the pushout of $S^{n-1}\hookrightarrow D^n$ and $f$.
\end{solution}

\begin{exercise}[Sphere quotient]\label{exr:viii05-ped-07}
Prove $D^n/S^{n-1}\cong S^n$.
\end{exercise}
\begin{hint}
Collapse the boundary to the point at infinity.
\end{hint}
\begin{solution}
Identifying the boundary of a closed ball to one point gives the one-point compactification of its interior $\mathbb R^n$, which is $S^n$.
\end{solution}

\begin{exercise}[Attachment functoriality]\label{exr:viii05-ped-08}
Explain why a homeomorphism of attaching data induces a homeomorphism of attached spaces.
\end{exercise}
\begin{hint}
Use the universal property of the quotient/pushout.
\end{hint}
\begin{solution}
Compatible homeomorphisms on the old space and disk descend to mutually inverse maps on the quotients.
\end{solution}

\begin{exercise}[Homotopic maps preview]\label{exr:viii05-ped-09}
Why should homotopic attaching maps often lead to the same homotopy type?
\end{exercise}
\begin{hint}
A homotopy provides a cylinder between the boundary maps.
\end{hint}
\begin{solution}
The homotopy allows one to interpolate the gluing in a mapping-cylinder construction; under standard CW/cofibration hypotheses this yields a homotopy equivalence of the adjunction spaces.
\end{solution}

\subsection*{Counterexamples and hypothesis tests}

\begin{exercise}[Dimension always increases]\label{exr:viii05-ped-10}
Test: attaching an $n$-cell always changes homotopy type.
\end{exercise}
\begin{hint}
An attachment can kill or create topology, but may also be redundant.
\end{hint}
\begin{solution}
False; depending on the attaching map and existing space, the resulting inclusion can sometimes be a homotopy equivalence or add contractible structure.
\end{solution}

\begin{exercise}[Attaching map interior]\label{exr:viii05-ped-11}
Test: an attaching map is defined on all of $D^n$.
\end{exercise}
\begin{hint}
Only the boundary is prescribed.
\end{hint}
\begin{solution}
False.  The attaching map has domain $S^{n-1}=\partial D^n$.
\end{solution}

\begin{exercise}[All attachments are wedges]\label{exr:viii05-ped-12}
Test: attaching a cell always gives a wedge with $S^n$.
\end{exercise}
\begin{hint}
That occurs for the constant attaching map under suitable based conditions.
\end{hint}
\begin{solution}
False; nontrivial attaching maps encode essential relations.
\end{solution}

\subsection*{Applications and investigations}

\begin{exercise}[Topological modeling]\label{exr:viii05-ped-13}
Why are cell attachments a useful finite description of spaces?
\end{exercise}
\begin{hint}
They separate local pieces from gluing data.
\end{hint}
\begin{solution}
A complicated space can be encoded by simple disks plus attaching maps, turning topology into finite combinatorial/algebraic data when only finitely many cells are used.
\end{solution}

\begin{exercise}[Relation encoding]\label{exr:viii05-ped-14}
Interpret a 2-cell attachment to a graph as imposing a relation on loops.
\end{exercise}
\begin{hint}
The attaching circle traces a word in the graph.
\end{hint}
\begin{solution}
The boundary loop becomes null-homotopic after filling it by the new disk, so its word is imposed as a relation in the fundamental group.
\end{solution}

\subsection*{Challenge problems}

\begin{exercise}[Moore-space preview]\label{exr:viii05-ped-15}
Analyze the cellular chain groups for one 0-cell, one 1-cell, and one 2-cell attached by degree $m$.
\end{exercise}
\begin{hint}
There is one generator in each dimension $0,1,2$.
\end{hint}
\begin{solution}
$C_2\cong C_1\cong C_0\cong\mathbb Z$, with $d_2$ multiplication by $m$ and $d_1=0$; hence $H_1\cong\mathbb Z/m$ for $m\ne0$ and $H_2=0$.
\end{solution}

\begin{exercise}[Attachment synthesis]\label{exr:viii05-ped-16}
Explain how the same attaching map participates in both homotopy and homology calculations.
\end{exercise}
\begin{hint}
It is a boundary map geometrically and induces algebraic boundary coefficients.
\end{hint}
\begin{solution}
Geometrically it determines which boundary sphere is glued to the lower skeleton; algebraically its degrees against lower cells become entries of the cellular differential, connecting homotopy data to homology.
\end{solution}

% END VOL08-EXPANSION VIII05
